\chapter{Переріз Пуанкаре й обернена задача}\label{ch:g07}

Це найдовша історія в проєкті. Спершу ми двічі помилилися в канонічних змінних (V-D3, R4) і з'ясували, що симплектична нейромережа~--- не наша ідея (V-D1, R3). Потім пілот D3 показав, що одномодова обернена задача тривіальна, а наше багатомодове узагальнення зламане (R8). Потім стенд \texttt{tokamak-3d-viz} зняв блокер і з'ясував, що в лінеаризованій постановці складність задачі задає не фізика, а вибір спостережуваного (VE21--VE25, VE30, VE31). Нарешті 2026-10-01 нелінійна перевірка за критеріями, записаними до запуску, закрила гіпотезу C5: базовий фіт не дійшов до рівня шуму на жодному стенді (R13), бо ландшафт нев'язки «скляний» (E40). Гамільтонів опис силових ліній і острови див. у~\cite{Manual}, розд.~2.

\section{Фізика: силова лінія як гамільтонова система}\label{sec:g07-ham}

У магнітних координатах $(\psit,\thetastar,\tor)$ поле токамака (з точністю до знака конвенції) має вигляд
\begin{equation}\label{eq:g07-B}
  \Bvec=\grad\psit\times\grad\thetastar+\grad\tor\times\grad\psi ,
\end{equation}
де $\psit$~--- тороїдальний потік на радіан, а $\psi$~--- полоїдальний. Уздовж силової лінії маємо
\begin{equation}\label{eq:g07-hamilton}
  \frac{\dd\thetastar}{\dd\tor}=\frac{\pa\psi}{\pa\psit}=\frac1q,\qquad
  \frac{\dd\psit}{\dd\tor}=-\frac{\pa\psi}{\pa\thetastar}.
\end{equation}
Це рівняння Гамільтона з півтора ступенями вільності: час~--- $\tor$, координата~--- $\thetastar$, імпульс~--- $\psit$, гамільтоніан $H=\psi$.

\begin{established}{E18--E20 (прочитано)}
\textbf{E18.} В осесиметрії $\psi=\psi(\psit)$ не залежить від $\thetastar$ і $\tor$. Тоді $\psit$ зберігається, система інтегровна, і переріз Пуанкаре дає \emph{рівно} контури $\psi$.
\textbf{E19.} Канонічний словник: час~--- $\tor$, координата~--- $\theta$, імпульс~--- \emph{тороїдальний} потік, гамільтоніан~--- \emph{полоїдальний}.
\textbf{E20.} Одна резонансна мода не дає хаосу: гамільтоніан $\psi_0(\psit)+A\cos(m\thetastar-n\tor)$ у гвинтовому куті $\xi=m\thetastar-n\tor$ не залежить від часу, тож система інтегровна~\cite{RegEst,Escande2024}.
\end{established}

\begin{established}{E39 (прочитано й перевірено, 2026-10-01)}
Обернене до E18 («немає осесиметрії~--- немає вкладених поверхонь») хибне. Landreman~\cite{Landreman2026} дав явні гладкі тороїдальні 3D-рівноваги МГД з $\grad p\ne0$ скрізь, крім осі, з вкладеними поверхнями і відхиленням від осесиметрії порядку одиниці: сімейство з $\iota\equiv2$ (усі лінії замкнені) і сімейство з широм. Ми перевірили статтю власним кодом: рівняння виконуються до $10^{-10}$--$10^{-12}$, числа $\beta_V=2/57$, $\iota(0)=2{,}28690$ відтворено, помилок не знайдено. \emph{Застереження, якого в статті немає:} у сімейства $\iota\equiv2$ саме поле $\Bvec$ має точну неперервну \emph{неізометричну} симетрію (еліптичне обертання, бо це лінійний образ поля Соловйова). Її порушують $|\Bvec|$, густина струму й тиск, тож як контрприклад до Ґреда це сімейство слабке. У сімейства зі широм афінної симетрії немає, але шир у прикладі статті лише $+0{,}04\,\%$~\cite{RegEst}. Код перевірки: \texttt{Our try/99-bibliography/landreman2026/}.  % lint-ok: E39 стосується неосесиметричних рівноваг, ι тут --- обертальне перетворення статті
\end{established}

\begin{refuted}{R4 (V-D3)}
«$\psi$ як канонічний імпульс». У нашому чернетковому формулюванні було саме так. Перевірка за~\cite{Escande2024} і Kallinikos та ін.~(2023; повних бібліоданих цієї праці в корпусі немає, лише згадка в реєстрі E19) виправила це до того, як помилка потрапила в матеріали хакатону~\cite{RegVerif}. Проте у вересні 2026 вона \emph{повторилася} в коді \texttt{tokamak-3d-viz}: ширину острова рахували з $\psi$ як імпульсом і завищували в $\sqrt q$ разів (N1, VR15). Виправлено 2026-09-29, деталі в розд.~\ref{ch:g08}. Ще один слід тієї самої помилки був у рядках 6--7 docstring \texttt{topology\_probe.py} (розд.~\ref{ch:g05}); його виправлено 2026-09-29.
\end{refuted}

\paragraph{Симплектичність.} Відображення Пуанкаре $(\psit,\thetastar)\mapsto(\psit',{\thetastar}')$ за один тороїдальний оберт є відображенням за час $2\pi$, породженим потоком гамільтонової системи. Тому воно зберігає площу (теорема Ліувілля): $\det\mathsf J=1$, де $\mathsf J$~--- якобіан відображення. Фізично це збереження магнітного потоку, $\grad\cdot\Bvec=0$. Мапа з $\det\mathsf J\ne1$ не є відображенням силових ліній, і все, що на ній виміряно, фізично нічого не означає.

\begin{refuted}{R3 (V-D1)}
«Симплектична нейромережа для відображень силових ліній~--- наша ідея». Спростовано: HénonNet~\cite{Burby2021} вчить відображення Пуанкаре силових ліній, і кожен його шар симплектичний за побудовою. Журнал верифікації додає ще одне: симплектичність є властивістю \emph{відображення}, а не статичного растра $\psi(R,Z)$. На осесиметричній сітці $65\times65$ вона вироджується в умову «$\psi$~--- добрий інваріант», тобто зводиться до GS-членів розд.~\ref{ch:g04}~\cite{RegVerif}.
\end{refuted}

\section{C5 і «Рівень 3»}\label{sec:g07-c5}

\begin{conjecture}{C5 (так її сформулювали; закрито 2026-10-01 $\to$ R13)}
Обернена задача Рівня~3 (спектр $(m,n)$) трактабельна за 48~год зі starter kit~\cite{RegConj}. Як гіпотезу перевірили й спростували~--- \S\ref{sec:g07-c5test}.
\end{conjecture}

\paragraph{Означення, додане ретроспективно 2026-09-29.} У корпусі «Рівень~3» вживався без означення, а Рівні~1--2 не визначені ніде. Реєстр гіпотез дає таке означення, прямо позначене як відновлене з контексту, а не цитоване~\cite{RegConj}:

\begin{quote}\itshape
Рівень~3~--- обернена задача на даних перетину Пуанкаре: за спостережуваною структурою силових ліній (проколи/екскурсії або слід на диверторі) відновити спектр збурення $(m,n)$, тобто амплітуди й фази мод.
\end{quote}

\noindent Про Рівні~1--2 нічого не стверджується.

\section{Пілот D3: Tokamap (E30)}\label{sec:g07-tokamap}

Перш ніж будувати інфраструктуру, пілот~\cite{g07D3} мав відповісти, чи обернена задача взагалі коректно поставлена. Стендом стала Tokamap~\cite{Balescu1998}: явне двовимірне відображення за один тороїдальний оберт, де $\Psi\ge0$~--- аналог тороїдального потоку (імпульс), а $T$~--- полоїдальний кут у частках $2\pi$.

\code{ourtry/03-deep-dives/D3-poincare-inverse/tokamap.py}{27}{54}{(профіль $q$, крок мапи, орбіта, $\det\mathsf J$)}

\begin{itemize}
\item \textbf{Рядки 27--28}: профіль $q(\Psi)=1+3\Psi^2$~--- \emph{наш} вибір, а не Балеску (docstring, рядки 14--16). Від нього залежить кожне кількісне число, але не якісна картина.
\item \textbf{Рядки 31--35} (\texttt{step}): $P=\Psi-1-\frac{x_L}{2\pi}\sin2\pi T$, $\Psi'=\tfrac12\big(P+\sqrt{P^2+4\Psi}\big)$, $T'=T+1/q(\Psi')-\frac{x_L}{4\pi^2}\frac{\cos2\pi T}{(1+\Psi')^2}\pmod 1$. Корінь квадратного рівняння гарантує $\Psi'\ge0$: вісь $\Psi=0$ мапа поважає автоматично, на відміну від стандартної мапи.
\item \textbf{Рядки 38--46} (\texttt{orbit}): ітерація з виходом при $\Psi<0$, $\Psi>20$ або нескінченності.
\item \textbf{Рядки 49--54} (\texttt{jacobian\_det}): $\det\mathsf J$ правими різницями з кроком $h=10^{-7}$. Похибка $\sim h$ задає «підлогу» $\sim10^{-8}$, нижче якої перевірка не бачить.
\end{itemize}

\begin{derivation}[чому Tokamap симплектична: твірна функція]
З рядка 33 випливає неявна форма $\Psi=\Psi'+\frac{x_L}{2\pi}\sin(2\pi T)\,f(\Psi')$, де $f(\Psi)=\Psi/(1+\Psi)$. Візьмемо твірну функцію другого роду
\[
  F(T,\Psi')=T\Psi'+G(\Psi')+h(T)\,f(\Psi'),\qquad G'=\frac1q,\qquad h(T)=-\frac{x_L}{4\pi^2}\cos2\pi T .
\]
Тоді $\Psi=\pa F/\pa T=\Psi'+h'(T)f(\Psi')$ і $T'=\pa F/\pa\Psi'=T+1/q(\Psi')+h(T)f'(\Psi')$, де $f'=(1+\Psi)^{-2}$. Це рівно рядки 32--34. Будь-яке відображення, породжене твірною функцією, симплектичне \emph{тотожно}. Нижче це знадобиться для R8.
\end{derivation}

\begin{established}{E30}
Tokamap реалізовано, і він точно симплектичний: $\max|\det\mathsf J-1|=5{,}6\cdot10^{-8}$ при $x_L=1{,}0$. $\Psi\ge0$ зберігається, перехід до хаосу монотонний~\cite{RegEst}.
\end{established}

Повторний запуск 2026-09-29 (\texttt{tokamap.py}) дав $\max|\det\mathsf J-1|$: $3{,}39\cdot10^{-9}$, $1{,}10\cdot10^{-8}$, $2{,}66\cdot10^{-8}$ і $5{,}64\cdot10^{-8}$ при $x_L=0$; 0,2; 0,5; 1,0. $\Psi\ge0$ виконується при $x_L=0{,}5$; 1,0; 2,0. Частка орбіт із розкидом $>0{,}3$ дорівнює 0 при $x_L\le0{,}5$, 0,083 при 0,8, 0,250 при 1,2 і 0,417 при 2,0 (рис.~\ref{fig:g07-tokamap}).

\begin{figure}[htbp]\centering
\includegraphics[width=\textwidth]{g07_tokamap.jpg}
\caption[Фазові портрети Tokamap]{Фазові портрети Tokamap ($q=1+3\Psi^2$) за 1\,500 ітерацій з 60 початкових умов. При $x_L=0$ маємо інваріантні кола (інтегровний випадок, аналог E18). Зі зростанням $x_L$ з'являються ланцюги островів на раціональних $q$ і розсіяні орбіти між ними. Власний розрахунок функцією \texttt{orbit} з \texttt{tokamap.py}; рисунок~--- \texttt{guide/figures/g07\_fig.py}.}\label{fig:g07-tokamap}
\end{figure}

\section{Пілот D3: спостережуване й ідентифіковність (E31, E32)}\label{sec:g07-ident}

\code{ourtry/03-deep-dives/D3-poincare-inverse/identifiability.py}{28}{46}{(усереднений профіль екскурсії)}

Спостережуване~--- радіальний профіль розкиду $\Psi$ за орбіту (\texttt{np.ptp} після 120 ітерацій прогріву), 32 радіуси. Ключовий рядок 34: для кожного радіуса беремо \texttt{N\_PHASE}$=24$ рівномірно розставлених початкових фаз $T_0$ зі спільним випадковим зсувом і усереднюємо. Перша версія брала \emph{одну} випадкову фазу.

\begin{established}{E31, E32}
\textbf{E31.} Спостережувана з однієї початкової фази непридатна: її шумова підлога 0,3244, а з усередненням по 24 фазах~--- 0,0006, у 540 разів краще. (Застереження: дві підлоги з різних стендів~--- зареєстровано 2026-09-29, див. нижче.)
\textbf{E32.} Відновити одну амплітуду тривіально: зміна $x_L$ на 3\,\% дає SNR 57, на 10\,\%~--- SNR 171. Для хакатону це не задача~\cite{RegEst}.
\end{established}

Повторний запуск: \texttt{identifiability.py} дає підлогу 0,0006 (максимум 0,0012), для +3\,\%~--- відносну зміну 0,0341 і SNR 56,6, для +10\,\%~--- 0,1028 і SNR 170,6. \texttt{multimode.py} відтворив підлогу 0,3244. \emph{Методичне застереження:} ці дві підлоги виміряно на \emph{різних} стендах. 0,3244 отримано на тримодовій мапі \texttt{multimode.py} ($a=0{,}3$, одна фаза), яка ще й не симплектична (R8), а 0,0006~--- на одномодовій Tokamap з $x_L=0{,}6$ і 24 фазами. Отже, «540$\times$» змішує ефект усереднення зі зміною стенда. Чистий ефект усереднення можна виміряти, поставивши \texttt{N\_PHASE}$=1$ в \texttt{identifiability.py} (практикум). Висновок E31, що одна фаза непридатна, це не скасовує: у \texttt{multimode.py} усі «сигнали» (0,38--0,43) лежать у межах шуму. Зареєстровано 2026-09-29 як застереження до E31: висновок про усереднення стоїть, відношення «540$\times$» кваліфіковано~\cite{RegEst,RegLog}.

\section{Пілот D3: зламане узагальнення (R8)}\label{sec:g07-r8}

\code{ourtry/03-deep-dives/D3-poincare-inverse/multimode.py}{30}{49}{(\emph{навмисно зламана} багатомодова мапа)}

Рядки 30--35: $V(T)=\sum_k\frac{a_k}{2\pi}\sin(2\pi m_kT+\varphi_k)$ і її похідна $V'(T)=\sum_ka_km_k\cos(\dots)$. Рядок 39 підставляє $V$ туди, де в Tokamap стоїть $\frac{x_L}{2\pi}\sin2\pi T$, а рядок 41 підставляє $V'/4\pi^2$ туди, де стоїть $\frac{x_L}{4\pi^2}\cos2\pi T$. Рядки 45--49 перевіряють $\det\mathsf J$ так само, як у \texttt{tokamap.py}. Симплектичність тут «перевіряється, а не припускається» (docstring), і перевірка її провалила.

\begin{refuted}{R8}
«Структуру Tokamap можна перенести на суму мод підстановкою $V'(T)$». Симплектичність ламається: $\max|\det\mathsf J-1|=1{,}7$ при $a=0{,}5$. Твірна функція такої підстановки не витримує~\cite{RegEst}. Повторний запуск дав $3{,}39\cdot10^{-9}$, $0{,}501$ і $1{,}69$ при $a=0$; 0,2; 0,5. Файл \texttt{multimode.py} лишено в репозиторії разом із провалом навмисно: «це найдешевший спосіб не повторити ту саму помилку»~\cite{g07D3}.
\end{refuted}

\begin{derivation}[де саме помилка (перевірка для путівника)]
За виведенням у \S\ref{sec:g07-tokamap}, у доданку $\Psi$ стоїть $h'(T)$, а в доданку $T'$~--- сама $h(T)$, тобто \emph{первісна}. Для суми мод $h'(T)=V(T)$, отже
\[
  h(T)=-\sum_k\frac{a_k}{4\pi^2m_k}\cos(2\pi m_kT+\varphi_k),
\]
і коефіцієнт має бути $a_k/m_k$, а не $a_km_k$, як дає $V'$. Для $m=1$ обидва збігаються, тому в одномодовій Tokamap підстановка «працює». Для $m_k=2,3,4$ вона помиляється в $m_k^2$ разів. Ми перевірили це чисельно тими самими точками й кроком, що в \texttt{multimode.py} (\texttt{guide/figures/g07\_multimode\_check.py}): з $h(T)$ замість $V'$ маємо $\max|\det\mathsf J-1|=9{,}5\cdot10^{-8}$ при $a=0{,}2$ і $3{,}6\cdot10^{-7}$ при $a=0{,}5$, тобто на рівні похибки різниць. Отже, це «шлях (1)» з рекомендацій D3 («вивести правильну твірну функцію»). Зареєстровано 2026-09-29 як \textbf{E35} (корінь R8; примітки в R8 і C5); скрипт~--- \texttt{Our try/03-deep-dives/D3-poincare-inverse/multimode\_fixed\_check.py}, а \texttt{multimode.py} лишається незмінним доказом R8~\cite{RegEst,RegLog}.
\end{derivation}

Пілот рекомендував інший шлях: взяти \texttt{pyoculus} (MIT), бо «власне узагальнення зламалось мовчки, і зламалось би так само вдруге». Блокер C5 сформульовано так: «потрібен валідний багатомодовий симплектичний стенд».

\section{Стенд без відображення: \texorpdfstring{$\Bvec=\grad\times(\psi_{\mathrm{total}}\grad\tor)$}{B = rot(ψ∇φ)} (VE21)}\label{sec:g07-stand}

\texttt{tokamak-3d-viz} обійшов проблему зовсім: відображення не будується, а поле задається ротором
\begin{equation}\label{eq:g07-curl}
  \Bvec=\grad\times(\psi_{\mathrm{total}}\grad\tor)+F\grad\tor,\qquad
  \psi_{\mathrm{total}}=\psi(R,Z)+\sum_kA_kg_k(\psiN)\cos(m_k\thetastar-n_k\tor+\alpha_k).
\end{equation}
Дивергенція ротора дорівнює нулю, тож $\grad\cdot\Bvec=0$ \emph{точно}, для будь-якого числа мод і будь-яких амплітуд, структурно, а не з точністю до похибки. Силові лінії далі просто інтегруються.

\code{viz/src/tokviz/perturbation.py}{114}{139}{(клас \texttt{Perturbation}: сума мод аналітично)}

Рядки 117--120 (docstring): нічого не растеризується на сітку EFIT. Кут $\thetastar$ береться з \texttt{ThetaStarField}, обвідна задана в замкненій формі, тож мода $m$ збуджує лише свій резонанс $q=m/n$. Це урок VE6: растеризація $\cos m\thetastar$ на $65\times65$ давала 2,4\,\% паразитної $m=1$ і фальшивий острів. Рядок 130: амплітуди задаються як частка повного діапазону потоку $|\psib-\psiax|$. Рядки 133--139 (\texttt{dpsi}) обчислюють суму $A_kg_k\cos(m_k\thetastar-n_k\tor+\alpha_k)$ з \eqref{eq:g07-curl}.

\code{viz/src/tokviz/fieldline.py}{51}{60}{(\texttt{FieldLines.B}: компоненти поля)}

Рядок 56: $\BR=-R^{-1}\pa_Z\psi$, $\BZ=R^{-1}\pa_R\psi$, $\Btor=F/R$ (\cite{Manual}, розд.~3). Рядки 57--59 додають поле збурення від підключеного «бекенда». Для \texttt{Perturbation} це $\delta\BR=-R^{-1}\pa_Z\delta\psi$, $\delta\BZ=R^{-1}\pa_R\delta\psi$, $\delta\Btor=0$ (\texttt{delta\_B}, рядки 160--170 у \texttt{perturbation.py}).

\code{viz/src/tokviz/fieldline.py}{70}{95}{(\texttt{\_rhs}: рівняння силової лінії з $\tor$ як часом)}

Рядки 71--76: стан~--- вектор $(R_1..R_n,Z_1..Z_n)$ для всієї пачки ліній, а лінії поза сіткою позначаються як «мертві». Рядки 78--83 обчислюють те саме поле, що в \texttt{B}. Рядки 85--90 дають $\dd R/\dd\tor=R\BR/\Btor$ і $\dd Z/\dd\tor=R\BZ/\Btor$. Коментар пояснює, чому рівняння записано через $\Bvec$, а не через $\psi$: так бекенд із \emph{тороїдальним} збуренням (реальні котушки, розд.~\ref{ch:g09}) теж працює. Рядки 91--94 заморожують лінії, що вийшли за межу, щоб одна втікачка не зламала спільний адаптивний крок усієї пачки. Інтегрує DOP853 (\texttt{trace}, рядки 101--134).

\begin{established}{VE21}
Блокер батьківської гіпотези C5 знято структурно. Відображення не використовується, $\Bvec=\grad\times(\psi_{\mathrm{total}}\grad\tor)$ має нульову дивергенцію точно для будь-якого числа мод; перевірено з трьома одночасними модами 2/1, 5/2, 3/1. \emph{Застереження:} поле зберігає потік точно, а інтегратор DOP853~--- ні: його дрейф $10^{-7}$ за 200 обертів (VE5)~\cite{RegViz}.
\end{established}

Ще одне чесне застереження з docstring \texttt{perturbation.py} (рядки 35--40): це \emph{заданий}, вакуумоподібний збурювач, а не самоузгоджений відгук плазми на котушки. Вакуумне поле систематично \emph{завищує} стохастичність, бо плазма екранує резонансні компоненти.

\section{Ідентифіковність як обумовленість}\label{sec:g07-cond}

\begin{outofcorpus}[лінеаризована обернена задача]
Лінеаризуємо спостережуване навколо точки $\vect p_0$ (амплітуди й фази): $\delta\vect y\approx\mathsf J\,\delta\vect p$. Після знерозмірення (відносна зміна $\vect y$ на відносну зміну амплітуди і на опорну зміну фази $\pi/4$) робимо SVD: $\mathsf J_n=\mathsf U\,\mathrm{diag}(s_1\ge\dots\ge s_P)\,\mathsf V^{\mathsf T}$. Шум $\sigma$ у даних дає похибку оцінки вздовж $i$-го правого сингулярного вектора порядку $\sigma/s_i$. Звідси число обумовленості $\kappa=s_1/s_P$ і «найслабший напрямок» $s_P/s_1$~--- параметр, який визначено найгірше. Вибір одиниць у $\kappa$ завжди є рішенням, і його треба назвати.
\end{outofcorpus}

\code{viz/src/run_inverse.py}{52}{62}{(\texttt{build\_seeds}: засів на рівнях і фазах)}

Для кожного радіального рівня \texttt{s} і кожної з \texttt{n\_phase} рівновіддалених фаз $\thetastar$ функція \texttt{point\_at} знаходить точку $(R,Z)$ на поверхні. Масив \texttt{lvl} запам'ятовує, якому рівню належить кожне зерно.

\code{viz/src/run_inverse.py}{65}{86}{(\texttt{observable}: дві редукції тих самих трас)}

Рядки 68--73 записують амплітуди й фази в моди та будують \texttt{FieldLines}. Рядки 74--78: трасування на \texttt{n\_turns} обертів, екскурсія $\max\psiN-\min\psiN$ на орбіту, \texttt{nan} для ліній, що вийшли. Рядки 79--86~--- серце порівняння: \emph{усереднене} спостережуване (середнє по фазах на рівень, як у пілоті) і \emph{фазово-роздільне} (кожне зерно окремо), обчислені з \emph{тих самих} трас.

\code{viz/src/run_inverse.py}{192}{212}{(знерозмірення і контроль обсягу даних)}

Рядки 196--203 ділять $\mathsf J$ на $|\vect y|$ і множать на масштаби параметрів ($A_0$ для амплітуд, $\pi/4$ для фаз). Рядки 205--212~--- контроль конфаундингу. Роздільне спостережуване має у \texttt{n\_phase} разів більше чисел, тож його підвибірку до \emph{тієї самої} кількості рядків, що в усередненого, аналізують окремо. Далі (рядки 214--263) для кожної редукції йдуть SVD і синтетична інверсія: $\delta\vect p_{\mathrm{true}}\sim\mathcal N(0;0{,}25^2)$, шум 0--10\,\%, 40 спроб, МНК.

\section{Що вийшло: VE22--VE25}\label{sec:g07-ve}

Запуск \texttt{python src/run\_inverse.py} \texttt{--out data/inverse} \texttt{--modes "2/1,5/2,3/1"} 	exttt{--fit-phases} з 22 рівнями, 6 фазами і 28 обертами (DIII-D 203702, кадр 75, $\qnf=3{,}699$). Числа нижче звірено з \texttt{data/inverse/inverse.json}.

\begin{center}\small
\begin{tabular}{lcccc}
\toprule
спостережуване & рядків & $\kappa$ & найслабший & похибка фаз при 10\,\%\\
\midrule
усереднене & 22 & 17,9 & 0,056 & 94,6\,\%\\
роздільне, підвибірка & 22 & 8,2 & 0,122 & 52,1\,\%\\
роздільне, повне & 132 & 5,8 & 0,172 & 14,2\,\%\\
\bottomrule
\end{tabular}
\end{center}

\begin{established}{VE22--VE25}
\textbf{VE22.} Відновити самі амплітуди тривіально: для двох рознесених мод $\kappa=1{,}2$, розділення мод 98,8\,\%, похибка 1:1 із шумом. Це продовжує висновок пілота («занадто легко»), а не спростовує його.
\textbf{VE23.} Складність задає спостережуване, а не фізика: для трьох упакованих мод з амплітудами \emph{і} фазами $\kappa=17{,}9$ (усереднене) проти 5,8 (роздільне); похибка фаз при 10\,\% шуму 94,6\,\% проти 14,2\,\%.
\textbf{VE24.} Перевага роздільності ділиться надвоє. \emph{Фазова інформація} дає 2,2$\times$ в обумовленості і 1,8$\times$ в точності фаз, \emph{обсяг даних}~--- ще 1,4$\times$ і 3,7$\times$. Початкове порівняння 22 проти 132 рядків було конфаундоване.
\textbf{VE25.} Затиснута мода втрачає визначеність лише разом із фазовою інформацією: в усередненому спостережуваному найслабший амплітудний напрямок~--- $A[5/2]$ (sv 1,06 проти 3,86), у роздільному він піднімається до sv 8,88 з 12,55~\cite{RegViz,g07Inverse}.
\end{established}

Висновок INVERSE §3: задача трактабельна, а її складність~--- \emph{ручка дизайну} змагання. Амплітуди, рознесені моди й усереднений профіль дають розминку. Амплітуди й фази упакованих мод із шумом 3--10\,\% дають справжню задачу. Усереднений профіль разом із фазами дає нерозв'язну. Амплітуди без фаз для двох мод мають $\kappa=1{,}2$, з фазами~--- 11,8. \emph{Після R13 (\S\ref{sec:g07-c5test}) це висновок про лінеаризовану задачу:} $\kappa$ тут порахована скінченними різницями з кроком 15\,\% амплітуди й 0,3~рад, тобто це січна, що усереднює нерівності ландшафту, і локальний ландшафт вона не описує.

\section{Слід на стінці і число мод: VC6, VC7 \texorpdfstring{$\to$}{→} VE30, VE31}\label{sec:g07-campaign}

\code{viz/src/run_footprint.py}{105}{113}{(спостережуване сліду: три гладкі числа на орбіту)}

Зерна лежать у SOL, трохи за сепаратрисою (\texttt{sol\_seeds}), і трасуються до стінки (\texttt{trace\_to\_wall}). На кожне зерно спостережуване дає довжину зв'язку в обертах і позицію удару $(\cos2\pi s,\sin2\pi s)$. Сира гістограма ударів для скінченних різниць не годиться, бо лічильники комірок стрибають (VE32). Модельне рішення названо явно: загасання збурення за сепаратрисою розширено з 0,05 до 0,60 (\texttt{--taper}), інакше поле придушене саме там, де утворюється слід.

\begin{established}{VE30 (колишня VC6)}
Слід на стінці зважений до краю. Обумовленість на сліді 9,3, проти 4,1 для фазово-роздільної екскурсії і 22,1 для усередненої; похибка фаз при 10\,\%~--- 27,3\,\% проти 19,3\,\%. Слід на 40\,\% дешевший (126~с проти 210~с), бо лінії гинуть за 0,7--6,6 обертів. Найсильніший напрямок (sv 7,21) на 99,9\,\% складається з фази 3/1, найслабший (sv 0,77)~--- з фази 2/1, тобто вдев'ятеро гірше: 3/1 сидить на $\psiN=0{,}900$, ближче до стінки, ніж 2/1 на 0,752~\cite{RegViz,g07Campaign}.
\end{established}

\code{viz/src/run_mode_scan.py}{21}{27}{(набори мод для VC7)}

Скрипт п'ять разів запускає \texttt{run\_inverse.py} з \texttt{--fit-phases} на 16 рівнях, 5 фазах і 24 обертах, щоразу щільніше пакуючи резонанси.

\begin{center}\small
\begin{tabular}{clcccc}
\toprule
$K$ & моди & $\kappa$ (усер.) & $\kappa$ (розд.) & найслабший (розд.) & фази при 10\,\%\\
\midrule
2 & 2/1, 3/1 & 22,1 & 4,1 & 0,245 & 19,3\,\%\\
3 & +5/2 & 10,7 & 4,9 & 0,205 & 18,9\,\%\\
4 & +3/2 & 28,7 & 6,7 & 0,149 & 31,5\,\%\\
5 & +7/2 & 35,2 & 6,9 & 0,146 & 25,2\,\%\\
6 & +1/1 & 170,3 & 11,5 & 0,087 & 37,2\,\%\\
\bottomrule
\end{tabular}
\end{center}

\begin{established}{VE31 (колишня VC7)}
Обумовленість гіршає з числом мод плавно, без порогу. Від $K=2$ до $K=6$ масштабно-інваріантний найслабший напрямок монотонно падає з 0,245 до 0,087, а $\kappa$ роздільного спостережуваного росте з 4,1 до 11,5. У стрибку на $K=6$ є домішка артефакту: мода 1/1 на $\psiN=0{,}335$ майже вдвічі розтягує радіальний діапазон при незмінних 16 рівнях~\cite{RegViz,g07Campaign}.
\end{established}

\paragraph{Узгодженість чисел між прогонами.} Для тих самих трьох мод (2/1, 5/2, 3/1) усереднене спостережуване дає $\kappa=17{,}9$ в INVERSE (22 рівні, 6 фаз, 28 обертів) і 10,7 у скані (16 рівнів, 5 фаз, 24 оберти). Роздільне дає 5,8 і 4,9 відповідно. Числа не суперечать одне одному, бо налаштування різні. Проте вони показують, що $\kappa$ \emph{усередненого} спостережуваного чутлива до вибірки, тож абсолютні значення слід порівнювати лише в межах одного прогону. У скані обумовленість усередненого навіть не монотонна ($22{,}1\to10{,}7\to28{,}7$). Надійний індикатор тут~--- роздільне спостережуване.

\section{Перевірка C5 (2026-10-01): нелінійний фіт}\label{sec:g07-c5test}

Усе вище~--- \emph{лінеаризація}: обумовленість $\kappa$ порахована навколо однієї точки. До 2026-10-01 у C5 навіть не було спростовувача (у реєстрі стояв прочерк). Тому перевірку поставили так, як її розв'язувала б команда на хакатоні: базовий багатостартовий нелінійний МНК (\texttt{scipy.optimize.least\_squares}) відновлює 3 амплітуди і 3 фази з фазово-роздільної екскурсії при 3\,\% шуму. Критерії записали до запуску (\texttt{c5/THEORY.md}, SHA-256 \texttt{1ea23e22\ldots81e1a4})~\cite{RegConj,g07C5}:
\begin{itemize}
\item \emph{успіх}~--- усі амплітуди в межах 10\,\% і всі фази в межах $20^\circ$;
\item \emph{хибний мінімум}~--- старт поза допуском, але з $Q\le1{,}5\,Q_{\mathrm{true}}$, де $Q$~--- сума квадратів нормованих нев'язок (у звітах перевірки її позначено $\chi^2$, тут $\chi$ зайнята температуропровідністю), тобто інший розв'язок, що пояснює дані на рівні шуму;  % lint-ok: χ² згадано лише як назву зі звітів
\item \emph{стенд A}~--- виправлена багатомодова Tokamap (E35), моди $m=2,3,4$, 40 екземплярів по 16 стартів; \emph{стенд B}~--- \texttt{tokamak-3d-viz}, DIII-D 203702, моди 2/1, 5/2, 3/1, 3 екземпляри по 6 стартів;
\item \emph{спростовано}, якщо на будь-якому стенді успіхів менше за 50\,\% або хибних мінімумів понад 25\,\%. Бюджет «48~год» критерієм не є, час лише звітується.
\end{itemize}

\code{ourtry/03-deep-dives/D3-poincare-inverse/c5/stand_tokamap.py}{38}{46}{(спостережуване стенда A: екскурсія $\Psi$ за 400 ітерацій)}

\noindent\emph{Рядки 40--41:} старт~--- фіксована сітка зерен \texttt{PSI0}, \texttt{T0} (20 радіусів $\times$ 8 фаз, задана на початку файлу) без випадкового зсуву, тож пряма модель є детермінованою функцією параметрів; \texttt{lo} і \texttt{hi} накопичуватимуть мінімум і максимум. \emph{Рядки 42--45:} мапа ітерується 400 разів, і для кожного зерна оновлюються мінімум і максимум $\Psi$. \emph{Рядок 46:} спостережуване~--- розмах $\max-\min$, тобто екстремум по \emph{дискретних} точках орбіти. Саме цей рядок пояснить результат.

\code{ourtry/03-deep-dives/D3-poincare-inverse/c5/fit_common.py}{74}{92}{(один фіт: старт, МНК і критерії з \texttt{THEORY.md})}

\noindent\emph{Рядки 76--78:} спостереження з 3\,\% мультиплікативного шуму і $Q$ самої істини~--- еталон, з яким порівнюють знайдений мінімум. \emph{Рядки 79--81:} випадковий старт з апріорного розподілу (амплітуди лог-рівномірно, фази рівномірно) і фіт у змінних $(\ln a,\varphi)$ з якобіаном за скінченними різницями. \emph{Рядки 84--92:} допуски успіху (\texttt{within}) і означення хибного мінімуму (\texttt{spurious}), рівно як у критеріях.

\begin{center}\small
\begin{tabular}{@{}lcc@{}}
\toprule
 & стенд A (Tokamap, E35) & стенд B (\texttt{tokamak-3d-viz})\\
\midrule
успіхи & \textbf{0 з 40} & \textbf{0 з 3}\\
хибні мінімуми & 0\,\% & 0\,\%\\
найкращий $Q/Q_{\mathrm{true}}$ & 2,6--193, медіана 110 & 303, 484, 503\\
медіанна макс. похибка амплітуди / фази & 32\,\% / $77^\circ$ & 54\,\% / $144^\circ$\\
процесорний час на екземпляр & 16~с & 5,3--7,3~год\\
\bottomrule
\end{tabular}
\end{center}

\noindent При 10\,\% шуму (стенд A, лише для звіту) успіхів 2 з 40. Хибних мінімумів немає: жоден старт не знайшов \emph{іншого} розв'язку, що пояснює дані на рівні шуму, тож інформація в спостережуваному є. Але жоден старт не дійшов до рівня шуму \emph{взагалі}: розв'язувач зупиняється в $Q$, у десятки--сотні разів більшому за істинний~\cite{g07C5}.

\begin{refuted}{R13 (була гіпотеза C5)}
«Обернена задача Рівня~3 трактабельна за 48~год зі starter kit». Базовий нелінійний фіт: 0 з 40 успіхів на Tokamap і 0 з 3 на стенді \texttt{tokamak-3d-viz}. Спростовано для \emph{базового розв'язувача}, а не «нерозв'язно»: перешкода~--- оптимізація на негладкому ландшафті, а не ідентифіковність. VE21--VE25, VE30, VE31 лишаються вимірюваннями обумовленості локальної січної, але ландшафту не описують~\cite{RegEst}.
\end{refuted}

\begin{figure}[htbp]\centering
\includegraphics[width=\textwidth]{g07_c5_landscape.jpg}
\caption{Безшумова сума квадратів нев'язок $Q$ уздовж однієї фази $\varphi_1$ через істину (стенд A, екземпляр 0), логарифмічна шкала. Ліворуч: у межах $\pm20^\circ$ із кроком $0{,}5^\circ$ є 8 локальних мінімумів, і зсув на $0{,}5^\circ$ дає $Q=245$ в один бік і 2,9 в інший. Праворуч: повне коло, 24 мінімуми. Пунктир~--- рівень шуму $Q\approx N_{\mathrm{obs}}=160$. Власний розрахунок, код: \texttt{Our try/03-deep-dives/D3-poincare-inverse/c5/diag\_tokamap.py}, \texttt{guide/figures/g07\_c5\_landscape.py}.}
\label{fig:g07-c5}
\end{figure}

\begin{established}{E40 (post-hoc, на вердикт не впливала)}
Ландшафт нев'язки екскурсійного спостережуваного «скляний»: базовий фіт застрягає навіть поруч з істиною. Старт в істині лишається на місці ($Q/Q_{\mathrm{true}}=0{,}99$--$1{,}00$, тобто код коректний). Старт з істини $+5\,\%$ і $+10^\circ$ зупиняється при $Q/Q_{\mathrm{true}}=41$--$100$, а з $+20\,\%$ і $+45^\circ$~--- при 111--779. Механізм у рядку~46 \texttt{stand\_tokamap.py}: розмах~--- максимум по дискретних точках орбіти, тож малий зсув параметра змінює, \emph{яка} точка дає екстремум, і нев'язка стрибає (рис.~\ref{fig:g07-c5}). На стенді B скан пар фаз з кроком $30^\circ$ дав 8 і 14 локальних мінімумів, усі далеко від рівня шуму ($Q_0\ge2\,184\gg N_{\mathrm{obs}}=132$)~\cite{RegEst,g07C5}.
\end{established}

\paragraph{Що з цього випливає.} Задача Рівня~3 \emph{у цій постановці} (екскурсія як розмах і градієнтний фіт) не годиться як завдання базового рівня: команда без глобальної оптимізації чи сурогатної моделі не отримає нічого. Дві ідеї з \texttt{c5/README.md} \emph{не перевірені} і в реєстрі не стоять. Перша: висновок «складність задає спостережуване» може стосуватися лише згладженої (січної) задачі. Друга: гладше спостережуване (середнє чи квантиль замість розмаху, довші траси) може розгладити ландшафт і повернути задачу в межі базового фіту~\cite{g07C5}.

\section{Стан C5}\label{sec:g07-status}

\begin{itemcard}{R13 (була C5)}{закрито 2026-10-01: спростовано (реєстр) \quad·\quad \S\ref{sec:g07-c5test}}
Базовий нелінійний фіт не відновлює 3 амплітуди і 3 фази при 3\,\% шуму на жодному стенді (0 з 40, 0 з 3), хоча хибних мінімумів немає; ландшафт «скляний» (E40). Що лишилося від попередньої історії: VE21--VE25, VE30, VE31 як вимірювання локальної січної обумовленості; симплектичний багатомодовий стенд Tokamap (E35); залежність $\kappa$ від вибірки (17,9 в INVERSE проти 10,7 у скані мод для тих самих трьох мод). Межі, названі в INVERSE §4, теж лишаються: поле вакуумне, без відгуку плазми; один розряд і один кадр; слід на стінці без трасування многовидів~\cite{RegEst,RegConj,g07Inverse}.
\end{itemcard}

\begin{practicum}[обернена задача]
\textbf{Відтворити пілот} (з \texttt{fusion equilibrium challenge/starter}, $\approx$20~с, 40~с, 90~с): \texttt{.venv/bin/python "../../Our try/03-deep-dives/D3-poincare-inverse/tokamap.py"}, потім \texttt{multimode.py} і \texttt{identifiability.py}. \textbf{Стенд} (з \texttt{tokamak-3d-viz}, хвилини): \texttt{python src/run\_inverse.py --out data/inverse --quick}.

\textbf{Задача~1 (E31 чесно).} Поставте \texttt{N\_PHASE = 1} в \texttt{identifiability.py} і виміряйте підлогу на \emph{тій самій} одномодовій мапі. Яка частина «540$\times$» припадає на усереднення, а яка на зміну стенда?

\textbf{Задача~2 (R8 $\to$ вправа на валідацію).} Виведіть твірну функцію для суми мод (\S\ref{sec:g07-r8}) самостійно, виправте копію \texttt{multimode.py} (звірка: \texttt{multimode\_fixed\_check.py}, E35) і повторіть розділ «3. IDENTIFIABILITY» уже на симплектичній мапі. Чи лишаються спектри розрізненними над підлогою з 24 фазами?

\textbf{Перевірка C5} (з \texttt{fusion equilibrium challenge/starter}): \texttt{.venv/bin/python "../../Our try/03-deep-dives/D3-poincare-inverse/c5/run\_tokamap.py"} ($\approx$2~хв, 10 ядер) і \texttt{diag\_tokamap.py} ($\approx$1~хв); стенд B (\texttt{run\_viz.py}) іде $\approx$3~год.

\textbf{Задача~3 (після R13: гладке спостережуване).} Нелінійний фіт розмаху застрягає (\S\ref{sec:g07-c5test}). Замініть у копії \texttt{stand\_tokamap.py} розмах $\max-\min$ на гладшу величину (середнє $|\Psi-\Psi_0|$ або квантиль 90\,\% по орбіті) і повторіть \texttt{diag\_tokamap.py}. Чи зникають локальні мінімуми на розрізі по $\varphi_1$? Якщо так, запустіть \texttt{run\_tokamap.py} з новим спостережуваним. Це ще не гіпотеза реєстру: перед запуском запишіть критерії, як у \texttt{c5/THEORY.md}.
\end{practicum}
